\documentclass[11pt]{article}
\usepackage[utf8]{inputenc}
\usepackage[T1]{fontenc}
\usepackage{lmodern}
\usepackage{amsmath,amssymb,amsthm}
\usepackage{mathtools}
\usepackage{bm}
\usepackage{geometry}
\geometry{margin=1in}
\usepackage[numbers,sort&compress]{natbib}
\usepackage{hyperref}
\usepackage{xcolor}
\usepackage{cleveref}
\usepackage{booktabs}
\usepackage{microtype}
\emergencystretch=3em

% ---------- Theorem environments ----------
\theoremstyle{plain}
\newtheorem{theorem}{Theorem}[section]
\newtheorem{proposition}{Proposition}[section]
\newtheorem{lemma}{Lemma}[section]
\newtheorem{corollary}{Corollary}[section]
\newtheorem{conjecture}{Conjecture}[section]
\theoremstyle{definition}
\newtheorem{definition}{Definition}[section]
\newtheorem{remark}{Remark}[section]

% ---------- Custom commands ----------
\newcommand{\KL}[2]{D_{\mathrm{KL}}\mathopen{}\left(#1 \mathbin{\Vert} #2\right)\mathclose{}}
\newcommand{\Tr}{\operatorname{Tr}}
\newcommand{\tr}{\operatorname{tr}}
\newcommand{\E}{\mathbb{E}}
\newcommand{\R}{\mathbb{R}}
\newcommand{\GL}{\mathrm{GL}}
\newcommand{\SPD}{\mathrm{SPD}}
\newcommand{\Lie}{\mathfrak{gl}}
\newcommand{\so}{\mathfrak{so}}
\newcommand{\holonomy}{\operatorname{Hol}}
\newcommand{\softmax}{\operatorname{softmax}}
\newcommand{\diag}{\operatorname{diag}}
\newcommand{\dd}{\mathrm{d}}
\newcommand{\bmu}{\boldsymbol{\mu}}
\newcommand{\bSigma}{\boldsymbol{\Sigma}}
\newcommand{\bphi}{\boldsymbol{\phi}}
\newcommand{\bOmega}{\boldsymbol{\Omega}}
\newcommand{\bI}{\mathbf{I}}
\newcommand{\Fcal}{\mathcal{F}}
\newcommand{\Dcal}{\mathcal{D}}
\newcommand{\Hcal}{\mathcal{H}}
\newcommand{\Mcal}{\mathcal{M}}
\newcommand{\Ncal}{\mathcal{N}}
\newcommand{\Scal}{\mathcal{S}}
\newcommand{\Zcal}{\mathcal{Z}}
\newcommand{\Pcal}{\mathcal{P}}
\newcommand{\Gcal}{\mathcal{G}}
\newcommand{\Wmeta}{W}
\newcommand{\Smeta}{S_{\mathrm{meta}}}
\newcommand{\Tmeta}{T_{\mathrm{meta}}}
\newcommand{\bmeta}{\beta_{\mathrm{meta}}}
\newcommand{\Omegaij}{\Omega_{ij}}
\newcommand{\betaij}{\beta_{ij}}
\newcommand{\gammaij}{\gamma_{ij}}
\newcommand{\veps}{\varepsilon}

\title{Configurational Meta-Entropy of a Population of Variational Beliefs:\\
Toward a Thermodynamic Limit for the Gauge-Theoretic Free Energy Transformer}
\author{}
\date{}

\begin{document}
\maketitle

\begin{abstract}
The gauge-theoretic variational free energy (VFE) transformer assigns to a population
of agents a single scalar functional $\Fcal$ over their Gaussian belief tuples
$(\bmu_i,\bSigma_i,\bphi_i)$. A given value of $\Fcal$ can be realised by many distinct
belief configurations, which invites a statistical-mechanical reading: in the
large-population limit the log-count of belief configurations at fixed value of the
belief-sector free energy is a
candidate configurational meta-entropy $\Smeta(\Fcal_0)=\log\Wmeta(\Fcal_0)$, an
analogue of Boltzmann's $S=k\log W$. This note develops that construction and delimits
what can and cannot yet be claimed for it. We first separate three entropies that are
easily conflated: the intra-belief Gaussian entropy, the attention-distribution entropy
already present in $\Fcal$, and the configurational meta-entropy. We take the counting
measure on the belief sector to be the Riemannian volume form of the Fisher--Rao metric,
which Chentsov's theorem and its extension to parametric families fix up to a positive
scalar, the information-geometric counterpart of the Liouville measure, while noting
that the gauge sector requires a separately specified group measure together with a
quotient or regulator. We show that the coupling energy is rendered extensive by the
Kac normalisation implicit in row-stochastic softmax attention, provided the
attention-entropy contribution is absorbed into the row log-partition $-\tau\log Z_i$,
whose $1/N$ normalisation keeps each row bounded. We give the meta-entropy a large-deviation
form through a single declared reference measure and identify the McKean--Vlasov
stationary form of its constrained minimiser, keeping the directed row structure of attention rather than a symmetrised
pair potential. The conjugate meta-temperature is given an operational reading as the
noise level of stochastic natural-gradient belief updating through the Riemannian
Langevin equation, whose stationary law requires the metric-divergence drift together
with confinement. A solvable Gaussian instance is a linear-response calculation rather
than a delivered phase transition: its susceptibility diverges as the prior-anchoring
strength $\alpha\to0$, identifying the prior as the explicit field that breaks an
otherwise exact local $\GL^{+}(K)$ covariance of the interaction before regulation. We close with a
conjecture, conditional on an edge-local connection with nontrivial holonomy and on
extensive ground-state degeneracy, that geometric frustration leaves a residual
meta-entropy; in the vertex-frame transport actually used in the functional the loop
holonomy is identically trivial, so the conjecture concerns a distinct, edge-relaxed
regime.
\end{abstract}

\section{Motivation}
\label{sec:motivation}

The gauge-theoretic VFE transformer dispenses with learned projections and obtains all
of its representational capacity from iterative minimisation of a single free energy
functional over Gaussian belief tuples \citep{vaswani2017attention,friston2010free}. In
its canonical form the functional couples $N$ agents through
\begin{equation}
\label{eq:F}
\begin{aligned}
\Fcal
={}& \alpha\sum_{i}\KL{q_i}{p_i}
+ \lambda_h\sum_i \KL{s_i}{h}
+ \sum_{i,j}\Bigl[\betaij \KL{q_i}{\Omegaij q_j} + \tau \betaij\log\tfrac{\betaij}{\pi_{ij}}\Bigr]\\
&{}+ \sum_{i,j}\Bigl[\gammaij \KL{s_i}{\Omegaij s_j} + \tau \gammaij\log\tfrac{\gammaij}{\pi^{s}_{ij}}\Bigr]
- \sum_i\E_{q_i}[\log p(o_i\mid x_i)],
\end{aligned}
\end{equation}
with transport $\Omegaij = e^{\bphi_i}e^{-\bphi_j}$, congruence action
$\bSigma\mapsto\Omegaij\bSigma\Omegaij^{\top}$, effective temperature $\tau=\kappa\sqrt{K}$,
and attention weights obtained by minimising the row-simplex Lagrangian,
\begin{equation}
\label{eq:softmax}
\betaij=\softmax_j\bigl(\log\pi_{ij}-\KL{q_i}{\Omegaij q_j}/\tau\bigr),
\end{equation}
which reduces to $\softmax_j(-\KL{q_i}{\Omegaij q_j}/\tau)$ only for the uniform prior
$\pi_{ij}=1/N$ on a fixed row support; throughout we take $\tau>0$ and $\pi_{ij}>0$ with
$\sum_{j\in S_i}\pi_{ij}=1$ on that active support $S_i$ (the operative case being full all-to-all
support, $S_i=\{1,\dots,N\}$), so the row optimum is interior and the row-reduced functional is
smooth in the belief and frame variables. The transport in \cref{eq:F} is of vertex-frame
form, $\Omegaij=g_i g_j^{-1}$ with $g_i=e^{\bphi_i}$; this distinction will matter in
\cref{sec:frustration}, where an edge-local connection is required for nontrivial
holonomy. A single value of $\Fcal$ does not pin down the configuration
$\{(\bmu_i,\bSigma_i,\bphi_i)\}$: the beliefs can be distributed among the agents in many
ways at fixed cost. The count is taken over the belief and frame sector
$(\bmu_i,\bSigma_i,\bphi_i)$ at fixed M-step parameters $p_i,s_i,h$, so the counted
functional is the belief-sector free energy $\Fcal_{\mathrm{belief}}$ of
\cref{sec:measure} rather than the full $\Fcal$. This is the situation that, in physics, gives rise to entropy. The
question we address is whether the count of belief configurations compatible with a
given free energy defines a bona fide entropy for the theory, what measure it must be
taken against, whether it survives the limit of an astronomically large population, and
what new structure the gauge sector contributes that ordinary thermodynamics lacks.

We answer in a qualified and partly conjectural form, and the construction requires care
because $\Fcal$ is itself already a free energy rather than an energy, so the new object
is an entropy of a free-energy landscape. Throughout we distinguish what is exact and
verified, what is a standard result imported from statistical mechanics and information
geometry, and what is conjectural. \Cref{sec:three} separates three distinct entropies.
\Cref{sec:measure} fixes the counting measure on the belief sector and states what the
gauge sector additionally requires. \Cref{sec:limit} examines extensivity through the
row-reduced free energy. \Cref{sec:meta} gives the meta-entropy a large-deviation form
under a declared reference measure. \Cref{sec:temperature} gives the meta-temperature an
operational meaning. \Cref{sec:phase} works a solvable Gaussian response calculation.
\Cref{sec:gauge} treats the gauge sector and separates invariance from redundancy.
\Cref{sec:frustration} states the holonomy-frustration conjecture for the edge-local
regime. \Cref{sec:discussion} records scope and open problems.

\section{Three entropies at three levels}
\label{sec:three}

The largest source of confusion is that $\Fcal$ already contains entropy, so a
meta-entropy must be defined against what is already there. Three distinct objects must
be kept apart. The first is the intra-belief entropy of a single Gaussian,
$\Scal[q_i] = \tfrac12\log\det(2\pi e\bSigma_i)$, which is the term that makes each
agent's contribution a free energy rather than a bare energy. The second is the
attention-distribution entropy, the family of terms
$\tau\betaij\log(\betaij/\pi_{ij})$ in \cref{eq:F}; these measure the spread of one
agent's attention over its neighbours, are a within-configuration quantity, and are part
of the functional whose minimisation over $\betaij$ returns the softmax \cref{eq:softmax}; the model channel carries a structurally identical attention-distribution entropy in the model-coupling ($\gamma_{ij}$) terms of \cref{eq:F}, scoped out here together with the rest of the $\gamma$ block.
The third, which is the subject of this note, is the meta-entropy: the logarithm of the
number of belief configurations $\{(\bmu_i,\bSigma_i,\bphi_i)\}_{i=1}^{N}$ that realise a
prescribed value of $\Fcal$. The first two are properties of a single configuration; the
third is a property of the configuration space. Only the third is a new object, and only
the third is the analogue of Boltzmann's $S=k\log W$. Because $\Fcal$ is already a free
energy, the meta-entropy is an entropy of a free-energy landscape, a second-order
construction with the same formal status as the quenched free energies of disordered
systems \citep{mezard1987spin}.

\section{Microstate, macrostate, and the counting measure}
\label{sec:measure}

A microstate is a full configuration
$\omega=\{(\bmu_i,\bSigma_i,U_i)\}_{i=1}^{N}$ in the product manifold $\Mcal^{N}$, where
the single-agent state space is $\Mcal_1=\R^{K}\times\SPD(K)\times\Gcal$, with
$\GL^{+}(K)$ the ambient gauge group. The primitive counted coordinate of the gauge
sector is the Lie-algebra generator $\bphi_i$ restricted to the compact principal-log
chart $\{\bphi:\|\bphi\|_2\le\pi-\eta\}$ introduced in \cref{eq:container}; the frame
$U_i=e^{\bphi_i}$ is its parameterization and generates the transports. The chart image
is a strict subset of the exponential image, which for $K\ge2$ is in turn a strict subset of the
ambient group,
\begin{equation}
\label{eq:nesting}
\{e^{\bphi}:\|\bphi\|_2\le\pi-\eta\}\subsetneq\{e^{\bphi}\}\subsetneq\GL^{+}(K),
\end{equation}
the second inclusion being proper because the matrix exponential is not surjective onto
$\GL^{+}(K)$ for $K\ge2$: a real matrix has a real logarithm only if every Jordan block of
a negative eigenvalue occurs an even number of times \citep{culver1966existence}, so an
element such as $\diag(-2,-\tfrac12)\in\GL^{+}(2)$ has positive determinant yet no real
logarithm. We take $\Gcal=\GL^{+}(K)$, the identity component of $\GL(K)$ \citep{hall2015lie}, as the ambient group and count over the chart. The meta-entropy is defined at a fixed realisation of the slow M-step
sector. The fast E-step infers the belief and frame variables
$(\bmu_i,\bSigma_i,\bphi_i)$ while the priors $p_i$, the model channel $s_i$, the
hyper-prior $h$, and the model coupling are held fixed as the slow parameters that
the M-step adjusts on a separate timescale. Here \emph{quenched} names the
two-timescale procedure of disordered statistical mechanics, freezing the slow
realisation and counting only the fast degrees of freedom at fixed slow parameters
\citep{mezard1987spin}; it is not a claim that the frozen slow sector is itself
disordered. We make no average over the slow sector, so no annealed or replica step
is claimed. The large-deviation passage of \cref{sec:meta}, and the iid Sanov rate of
\cref{eq:sanov}, are stated in the operative homogeneous regime: the prior is common,
$p_i=p$, the observation potential is a common bounded per-agent functional, the slow
bounds are common, and the attention prior is the uniform $\pi_{ij}=1/N$ already
assumed. In this regime the frozen slow sector carries no disorder, the belief-sector
free energy depends on the configuration only through the empirical law of the fast
states, $\Fcal_{\mathrm{belief}}=N\Fcal[\rho_N]$, so $\rho_N$ is a sufficient
statistic, the fast population is exchangeable, and ordinary (non-disordered) iid
empirical-measure large deviations apply \citep[Thm.~6.2.10]{dembo1998large}. A
genuinely heterogeneous frozen slow sector (agent-dependent $p_i$ or observation
potentials) is the quenched-\emph{disorder} generalisation, in which
$\Fcal_{\mathrm{belief}}$ is no longer a function of the fast empirical law alone: it
requires augmenting the empirical state with the slow labels and a conditional, joint
fast--slow large-deviation rate \citep{comets1989large}, or an annealed average over a
declared distribution of the slow variables \citep{touchette2009large}; neither is
claimed here. The counted functional is
accordingly the belief-sector part of \cref{eq:F},
\begin{equation}
\label{eq:Fbelief}
\Fcal_{\mathrm{belief}}=\alpha\sum_i\KL{q_i}{p_i}
+\sum_{i,j}\Bigl[\betaij\KL{q_i}{\Omegaij q_j}+\tau\betaij\log\tfrac{\betaij}{\pi_{ij}}\Bigr]
-\sum_i\E_{q_i}[\log p(o_i\mid x_i)],
\end{equation}
the belief-sector restriction of $\Fcal$ that we count over $(\bmu_i,\bSigma_i,\bphi_i)$ at
fixed $p_i,s_i,h$. The hyper-prior term $\lambda_h\sum_i\KL{s_i}{h}$ depends only on
the fixed slow sector and is a configuration-independent additive offset that shifts
the zero of $\Fcal_0$ but not the count. The model-coupling block, by contrast, is
not a constant of the count: its transport $\Omegaij=e^{\bphi_i}e^{-\bphi_j}$ and the
resulting weight $\gammaij$ both depend on the counted frames $\bphi_i$, so it is a
slow-sector interaction that is excluded by scoping the count to
$\Fcal_{\mathrm{belief}}$, not by being configuration independent. A macrostate is the
value $\Fcal_0$ of $\Fcal_{\mathrm{belief}}$, together with any
further macroscopic constraints that play the role of a container volume. Throughout the
count, $\Fcal_{\mathrm{belief}}$ abbreviates the row-reduced functional
$\Fcal_{\mathrm{belief}}^{\mathrm{red}}(\omega)$ obtained by evaluating \cref{eq:Fbelief} at the
softmax stationary point \cref{eq:softmax}, so that the attention weights $\betaij$ are slaved to the
microstate through the row log-partition $-\tau\log Z_i$ of \cref{eq:rowfree} and
$\Fcal_{\mathrm{belief}}^{\mathrm{red}}$ is a function of $(\bmu_i,\bSigma_i,\bphi_i)$ alone; the
$\betaij$ are not independent microstate coordinates and are distinct from the canonical conjugate
$\bmeta$ introduced in \cref{sec:temperature}. The
microcanonical meta-entropy is
\begin{equation}
\label{eq:Smicro}
\Smeta(\Fcal_0)=\log\Wmeta(\Fcal_0),\qquad
\Wmeta(\Fcal_0)=\int_{\Dcal^{N}}\delta\bigl(\Fcal_{\mathrm{belief}}^{\mathrm{red}}(\omega)-\Fcal_0\bigr)\mu_1^{\otimes N}(\dd\omega).
\end{equation}
The Dirac density $\delta(\Fcal_{\mathrm{belief}}-\Fcal_0)$ is shorthand for the thin energy shell
$\Fcal_{\mathrm{belief},N}/N\in[f-\veps,f+\veps]$ at $f=\Fcal_0/N$, with the limit $N\to\infty$ taken first
and $\veps\to0$ afterwards; under a regular-value assumption on $\Fcal_{\mathrm{belief}}$ the coarea prefactors
that convert the shell volume into a surface density on the level set $\{\Fcal_{\mathrm{belief}}=\Fcal_0\}$ are
subexponential in $N$ and so leave the entropy rate $\Smeta/N$ unchanged
\citep[\S5.4]{touchette2009large} and \citep{ellis1985entropy}. This shell-volume to LDP-rate
equality is used at regular values of $f$; at critical values and at the ground value $f_0$ of
\cref{sec:frustration}, where the level set can fail to be regular and the entropy density can lose
smoothness, the equality is read through the thick-shell limit above rather than the surface coarea form.
In Hamiltonian statistical mechanics the counting measure $\dd\mu=\prod \dd q \dd p$ is
canonical because the dynamics preserves it (Liouville's theorem)
\citep{goldstein1980classical}. The VFE dynamics is gradient descent on $\Fcal$, not a
symplectic flow, so a different principle must select the measure. The information
geometry of the belief manifold supplies it on the statistical sector: by Chentsov's
theorem and its extension to parametric families the Fisher--Rao metric is, up to scale,
the unique Riemannian metric invariant under Markov morphisms
\citep{chentsov1982statistical,amari2000methods,nielsen2020elementary}. The natural
single-agent reference is therefore the product of the Fisher--Rao volume on the
statistical factor and a group measure on the gauge factor,
\begin{equation}
\label{eq:measure}
\dd\mu_1=\dd\mathrm{Vol}_{\mathrm{FR}}(\bmu,\bSigma)\times\dd\mu_{\mathrm{Haar}}^{\Gcal}(U),
\qquad U\in\Gcal,
\end{equation}
with the affine-invariant Fisher--Rao volume on $\R^{K}\times\SPD(K)$, understood as the
Riemannian volume form $\sqrt{\det g_{\mathrm{FR}}}$ of the Fisher--Rao metric that the
theorem fixes up to a positive scalar (so the volume is determined up to a multiplicative
constant $c$, which shifts $\Smeta$ by the configuration-independent additive term $N\log c$ and hence
the entropy density $\Smeta/N$ by $\log c$; the Chentsov scale therefore fixes the additive zero of the
entropy density, and only entropy differences and $f$-derivatives are scale-invariant unless the scale
convention is declared), and the
Haar measure on the group $\Gcal$ itself. Two cautions attach to the gauge factor and are
carried through the rest of the note. First, Chentsov's theorem fixes the statistical
metric but says nothing about the group measure, which must be specified separately;
Haar measure lives on the group, not directly on the Lie-algebra coordinates $\bphi$. In
exponential coordinates $U=e^{\bphi}$ the group measure pulls back to a Jacobian density
$J_{\exp}(\bphi)\dd\bphi$ on a chart, but for $\GL^{+}(K)$ the exponential map is not a
global diffeomorphism \citep{higham2008functions}, so the chart and its Jacobian must be
stated. Second, $\Gcal=\GL^{+}(K)$ is noncompact, so the total gauge volume diverges and
$\Wmeta(\Fcal_0)$ is finite only after a gauge slice or quotient, or a regulator such as
a bound $\|\bphi_i\|\le c$, is imposed; we return to this container requirement in
\cref{sec:gauge}. With these provisos the correspondence is exact in role on the
statistical sector and interpretive in content: Chentsov's theorem is to the belief part
of the meta-entropy what Liouville's theorem is to thermodynamic entropy, each singling
out a natural measure for counting states, the present one curved rather than flat.

\section{Extensivity through the row-reduced free energy}
\label{sec:limit}

For the meta-entropy to scale with population size, the counted functional
$\Fcal_{\mathrm{belief}}$ must be extensive, $\Fcal_{\mathrm{belief}}=O(N)$, and the route to
extensivity is more delicate than row normalisation alone. At fixed row-stochastic weights, $\sum_j\betaij=1$, the coupling energy
$\sum_{i,j}\betaij \KL{q_i}{\Omegaij q_j}$ is a sum of $N$ convex combinations, hence
$O(N)$ once each divergence is bounded by the container of \cref{sec:measure}. The
attention-entropy term, however, is not bounded by $O(N)$ on its own: with $\pi_{ij}=1/N$,
\begin{equation}
\tau\sum_{j}\betaij\log\frac{\betaij}{\pi_{ij}}
=\tau\bigl(\log N-H(\beta_{i\cdot})\bigr)\in[0,\tau\log N],
\end{equation}
per row, where $H$ is the Shannon entropy of the row, so the naive total is
$O(N\log N)$ and is configuration dependent through the row concentrations $H(\beta_{i\cdot})$.
The correct extensive object is the row-reduced free energy obtained by inserting the
softmax stationary point \cref{eq:softmax} back into \cref{eq:F}. The energy and entropy
of row $i$ then combine into a log-partition,
\begin{equation}
\label{eq:rowfree}
\min_{\beta_{i\cdot}}\Bigl[\sum_j\betaij\KL{q_i}{\Omegaij q_j}+\tau\sum_j\betaij\log\tfrac{\betaij}{\pi_{ij}}\Bigr]
=-\tau\log Z_i,\qquad
Z_i=\sum_j\pi_{ij} e^{-\KL{q_i}{\Omegaij q_j}/\tau}.
\end{equation}
With $\pi_{ij}=1/N$ this is
$-\tau\log Z_i=\tau\log N-\tau\operatorname{logsumexp}_j\bigl(-\KL{q_i}{\Omegaij q_j}/\tau\bigr)$.
The $\tau\log N$ does not separate off as a removable reference: it is the Kac
normalisation that keeps the row bounded. Equivalently one writes the reduced row
contribution with the $1/N$ inside the logarithm,
\begin{equation}
\label{eq:rowkac}
-\tau\log Z_i=-\tau\log\Bigl(\tfrac1N\sum_j e^{-\KL{q_i}{\Omegaij q_j}/\tau}\Bigr)\in[0,M],
\end{equation}
where $M$ is the supremum of the row divergence over the container of \cref{sec:measure}.
Each summand $e^{-\KL{q_i}{\Omegaij q_j}/\tau}\in[e^{-M/\tau},1]$, so the averaged
partition function $Z_i$ lies in $[e^{-M/\tau},1]$ and each $-\tau\log Z_i$ lies in the
finite range $[0,M]$. Subtracting $\tau\log N$ per row would instead leave
$-\tau\sum_i\operatorname{logsumexp}_j(-\KL{q_i}{\Omegaij q_j}/\tau)$, which is
$O(N\log N)$ because $\operatorname{logsumexp}_j(-\KL{q_i}{\Omegaij q_j}/\tau)=\log N+O(1)$;
the $1/N$ must be kept inside the logarithm. Summed with the $1/N$ retained, the reduced
free energy $\sum_i(-\tau\log Z_i)$ is therefore $O(N)$. This is the precise content of the
Kac prescription, which renders long-range interactions extensive by normalising the
coupling by the number of partners \citep{kac1963vanderwaals}, the construction underlying
the Curie--Weiss model and mean-field theory. Row-stochastic softmax attention is thus a
Kac-normalised all-to-all coupling, so the reduced functional is extensive. Boundedness of
each $-\tau\log Z_i$ gives the $O(N)$ scaling but not, on its own, convergence of the
free-energy density $\frac1N\sum_i(-\tau\log Z_i)$; existence of that limit is the
Kac/Laplace-principle result \citep{kac1963vanderwaals,ellis1985entropy} and is realised
in \cref{sec:meta}, where the density is a weakly continuous functional of the empirical
measure that concentrates under Sanov. A mean-field variational limit is then the
appropriate leading-order description under these hypotheses, provided one works with the
reduced functional with the $1/N$ retained inside the row log-partition.

\section{The meta-entropy and its large-deviation form}
\label{sec:meta}

We give the meta-entropy a large-deviation form through a single declared reference
measure, following the microcanonical route and recording the canonical alternative for
completeness. Let $\rho_N=\tfrac1N\sum_i\delta_{(\bmu_i,\bSigma_i,U_i)}$ be the empirical
belief distribution, and let the reference be the product law $\mu_1^{\otimes N}$ built
from the single-agent measure \cref{eq:measure} restricted to the single-agent compact
container $\Dcal$ of \cref{eq:container} and normalised there to a probability law, so
that $\mu_1^{\otimes N}$ lives on the product container $\Dcal^{N}$ and $\rho_N$ is a
probability law on the single-agent space $\Dcal$. This product reference over the frame
coordinates $U_i$ is the correct starting point under the active-frame interpretation of
\cref{sec:gauge}, in which the local transformations are genuine changes of epistemic
configuration; if instead the frame directions are true gauge redundancies, the empirical law
must be carried on a gauge slice and the reference rebuilt as the corresponding quotient or
Faddeev--Popov measure \citep{faddeev1967feynman} before Sanov is applied, since quotienting
removes extensive orbit volume and can alter whether the soft frame directions contribute to
$\Smeta$. By Sanov's
theorem the empirical measure satisfies a large-deviation principle with the relative
entropy as rate,
\begin{equation}
\label{eq:sanov}
\Pr(\rho_N\approx\rho)\asymp e^{-N\KL{\rho}{\mu_1}}.
\end{equation}
The macroscopic free energy is a functional of the population law, and we keep the
directed (row) structure of attention rather than symmetrising. Write the continuum
single-agent state as $z=(\bmu,\bSigma,U)$, with $q=\Ncal(\bmu,\bSigma)$ its Gaussian
part and transport $\Omega(z,z')=U U'^{-1}$ determined by the frame pair $(U,U')$, so that
the empirical law of \cref{eq:sanov} is a law on these states. In the mean-field limit
each query state $z$ sees the population as its source distribution, drawn from a source
prior $\pi(\dd z'\mid z)$ that is the continuum image of the finite attention prior
$\pi_{ij}$; its row free energy is the log-partition of \cref{eq:rowfree} with the sum
over partners replaced by an integral against $\pi(\cdot\mid z)$, and the reduced functional is
\begin{equation}
\label{eq:Frho}
\Fcal[\rho]=\alpha\int \KL{q}{p}\rho(\dd z)
-\tau\int \rho(\dd z)\log\int e^{-\KL{q}{\Omega(z,z') q'}/\tau}\pi(\dd z'\mid z)+\cdots,
\end{equation}
where the attention weight $\beta(z,z')=e^{-\KL{q}{\Omega(z,z') q'}/\tau}/Z(z)$ is
normalised over the source prior by the row partition $Z(z)$, the inner integral of
\cref{eq:Frho}. For the uniform finite prior $\pi_{ij}=1/N$ on a fixed row support the
continuum source prior is the empirical law itself, $\pi(\dd z'\mid z)=\rho(\dd z')$, the
case treated below; then $Z(z;\rho)$ depends on the whole law $\rho$ through the row
normaliser. A structured or non-uniform $\pi_{ij}$ (a sparse interaction graph, a positional mask, or a
non-exchangeable kernel) falls outside the empirical-measure contraction of \cref{eq:contraction}, which
counts only the exchangeable fast states: there $\Fcal_N/N$ is no longer a continuous functional of $\rho_N$
alone, and a converging extension would have to augment the empirical state with the additional structure that
the kernel resolves, a graphon limit of the interaction array or a marked or pair empirical measure carrying
the source labels, which we do not pursue here. The trailing $+\cdots$ in \cref{eq:Frho} is the population observation term $\int V(z)\rho(\dd z)$, the $\rho$-average of a common bounded per-agent potential $V(z)=-\E_q[\log p(o\mid x)]$ of \cref{eq:F}; under the homogeneous slow sector this is a single $\rho$-functional bounded on $\Dcal$, with per-agent average $O(1)$, so the summed observation block is $O(N)$ and contributes extensively. No model-sector term enters because the count is scoped to $\Fcal_{\mathrm{belief}}$ at fixed slow parameters. No symmetric factor of $1/2$ appears, because the coupling is
row-directed; a symmetric double integral would be correct only after the conductances
have been symmetrised, which is a separate modelling choice. In the operative uniform-prior
case $\pi(\dd z'\mid z)=\rho(\dd z')$, on the compact container
$\Dcal$ the Gaussian divergences $\KL{q}{p}$ and $\KL{q}{\Omega(z,z') q'}$ are jointly
continuous and uniformly bounded, so the kernel $e^{-\KL{q}{\Omega(z,z') q'}/\tau}$ is bounded
away from zero and $\rho\mapsto\Fcal[\rho]$ is weakly continuous on the Prokhorov-compact
space $\Pcal(\Dcal)$; this discharges the continuity hypothesis of the contraction
principle \citep[Thm.~4.2.1]{dembo1998large} applied to the Sanov rate
\citep[Thm.~6.2.10]{dembo1998large}, provided the priors and the model sector $s_i$ obey
the same two-sided bound and the likelihood $-\E_q[\log p(o\mid x)]$ is bounded on
$\Dcal$. Combining \cref{eq:sanov}
with the constraint $\Fcal[\rho]=f$ through the contraction principle gives the
microcanonical meta-entropy density, up to the configuration-independent reference and
container terms,
\begin{equation}
\label{eq:contraction}
s(f)=\lim_{N\to\infty}\frac1N\log\Wmeta(Nf)
=-\inf\bigl\{ \KL{\rho}{\mu_1} : \Fcal[\rho]=f \bigr\}+\text{const}.
\end{equation}
The meta-entropy at free-energy density $f$ is thus minus the cheapest relative-entropy
cost, against the declared reference, of any population belief-distribution consistent
with $f$.

\begin{remark}[Explicit Gaussian form of the Sanov rate]
\label{rem:gaussianrate}
The rate $\KL{\rho}{\mu_1}$ is a relative entropy between probability laws on the
single-agent manifold $\Mcal_1$, not the Gaussian divergence of \cref{lem:affine}. On the Gaussian sector, with the gauge coordinate held fixed or marginalised against its Haar
factor so that the frame carries no relative entropy, the reference $\mu_1$ of \cref{eq:measure}
reduces to the Fisher--Rao volume normalised on the container $\Dcal$, with density against the
Riemannian volume the constant
$\dd\mu_1/\dd\mathrm{Vol}_{\mathrm{FR}}=1/\mathrm{Vol}_{\mathrm{FR}}(\Dcal)$; the determinant
density $c_K(\det\bSigma)^{-(K+2)/2}$ displayed below is this Gaussian-sector factor of the full
$\mathrm{Vol}_{\mathrm{FR}}\times\mu_{\mathrm{Haar}}$ product, not the joint density. The
Gaussian-sector marginal rate then reduces exactly, for every $\rho\ll\mu_1$ on $\Dcal$ (the rate
is $+\infty$ otherwise), to a negative differential entropy
taken against $\sqrt{\det g_{\mathrm{FR}}}$,
\begin{equation}
\label{eq:rateentropy}
\KL{\rho}{\mu_1}=-H_{\mathrm{FR}}[\rho]+\log\mathrm{Vol}_{\mathrm{FR}}(\Dcal),
\qquad
H_{\mathrm{FR}}[\rho]=-\int_{\Dcal}\frac{\dd\rho}{\dd\mathrm{Vol}_{\mathrm{FR}}}\log\frac{\dd\rho}{\dd\mathrm{Vol}_{\mathrm{FR}}}\dd\mathrm{Vol}_{\mathrm{FR}};
\end{equation}
this relative entropy is nonnegative and vanishes only at $\rho=\mu_1$, and being unbounded
above as $\rho$ concentrates it is a coercive, lower-semicontinuous Sanov rate. On the
Gaussian sector the Fisher metric is block-diagonal in $(\bmu,\bSigma)$ with mean block
$\bSigma^{-1}$ and covariance block proportional to the affine-invariant SPD metric (half it, in the Fisher normalisation), so the volume element
is $\sqrt{\det g_{\mathrm{FR}}}=c_K(\det\bSigma)^{-(K+2)/2}$, the product of the mean-block
weight $(\det\bSigma)^{-1/2}$ and the SPD volume $(\det\bSigma)^{-(K+1)/2}$
\citep{skovgaard1984riemannian,amari2000methods}; block-diagonality of the metric does not
make the joint volume factor into independent $\bmu$ and $\bSigma$ measures, the
$\bmu$-volume carrying the $\bSigma$-dependent weight $(\det\bSigma)^{-1/2}$. The
affine-invariant geodesic distance of that covariance block,
$d_R(\bSigma,\bSigma_0)=\lVert\log(\bSigma_0^{-1/2}\bSigma\bSigma_0^{-1/2})\rVert_F=(\sum_i(\log\lambda_i)^2)^{1/2}$
with $\lambda_i$ the eigenvalues of $\bSigma_0^{-1}\bSigma$, is the closed-form geodesic of
the Gaussian covariance geometry and is invariant under the congruence
$\bSigma\mapsto A\bSigma A^{\top}$ \citep{pennec2006riemannian}; the same-mean Fisher--Rao
distance equals $d_F=d_R/\sqrt2$ \citep{pinele2020fisher}. This distance fixes the volume
form in \cref{eq:rateentropy} but is not itself the rate, which is the law-to-law relative
entropy of \cref{eq:rateentropy} rather than the point-to-point Gaussian divergence of
\cref{lem:affine}: that point divergence between two same-mean covariance points agrees
with $\tfrac12 d_F^2=\tfrac14 d_R^2$ only to second order, the Fisher metric being the
Hessian of the Gaussian divergence \citep{amari2000methods,nielsen2020elementary}, and the
metric-normalisation convention relating $d_F$ and $d_R$ shifts
$\log\mathrm{Vol}_{\mathrm{FR}}(\Dcal)$ by a configuration-independent constant as in the
Chentsov-scale remark of \cref{sec:measure}. The constrained minimiser
$\rho^{\star}=\arg\min\{\KL{\rho}{\mu_1}:\Fcal[\rho]=f\}$ then takes exponential-family form for
the fixed-covariance means-only sub-family, where $\Mcal_1=\R^{K}$ carries the constant
metric $\bSigma^{-1}$ and the entropy-maximising minimiser under a quadratic $\Fcal$ is the
$I$-projection of the uniform carrier $\mathbf 1_{\{\lVert\bmu\rVert\le R\}}\dd\mathrm{Vol}$ onto the
moment family it constrains, namely the truncated exponential tilt
$\rho^{\star}(\bmu)\propto e^{\mathbf a\cdot\bmu-\tfrac12\bmu^{\top}\mathbf A\bmu}\mathbf 1_{\{\lVert\bmu\rVert\le R\}}$
with carrier the ball, where $\mathbf a\in\R^{K}$ and the symmetric $\mathbf A$ are the conjugate
parameters fixed by $\Fcal[\rho^{\star}]=f$ \citep{csiszar1975iprojection,cover2006elements}. The
all-to-all means-only $\Fcal[\rho]$ is nonlinear in $\rho$ through an interaction term
$-J\lVert\E\bmu\rVert^{2}$ of the form realised in \cref{prop:gaussian}, so the linear coefficient
$\mathbf a$ vanishes and $\mathbf A$ collapses to the scalar form $\mathbf A=\theta\bI$ only in the
centered, rotationally symmetric, zero-field subcase, the truncated-Gaussian family with sufficient
statistic $\lVert\bmu\rVert^{2}$ and a single multiplier $\theta\in\R$; on the compact ball $\theta$
ranges over all of $\R$, with $\theta>0$ selecting the low-second-moment branch and $\theta<0$ the
high-second-moment branch concentrated near the boundary, $\theta>0$ being recovered as the operative
sign only in the boundary-negligible limit below. Its exact rate is
$\KL{\rho^{\star}}{\mu_1}=-H_{\mathrm{FR}}(\rho^{\star})+\log\mathrm{Vol}(\{\lVert\bmu\rVert\le R\})$, and only in the
boundary-negligible limit, where the Gaussian width $\theta^{-1/2}$ is small relative to $R$ so the
truncation is inactive and $\rho^{\star}$ tends to the full Gaussian $\Ncal(0,S^{\star})$ with
$S^{\star}=\theta^{-1}\bI$ of \cref{prop:gaussian}, does this reduce to the displayed elementary value
$\KL{\rho^{\star}}{\mu_1}=-\tfrac12\log\det(2\pi e S^{\star})+\log\mathrm{Vol}(\{\lVert\bmu\rVert\le R\})$. With the covariance sector active the minimiser is implicit rather than elementary: its functional
derivative carries nonlinear covariance terms from the Gaussian divergences $\KL{q}{p}$ and
$\KL{q}{\Omega q'}$ and from the row normalisers of \cref{eq:Frho}, so the stationarity condition
does not solve in closed form and the residual $\rho^{\star}$ stays a McKean--Vlasov self-consistency
problem recorded as open. The geometry sharpens the same warning: the two-point Fisher--Rao distance
on the joint $(\bmu,\bSigma)$ manifold itself has no known closed form
\citep{nielsen2020elementary,pinele2020fisher}.
\end{remark} Stationary points of this constrained variational problem satisfy a self-consistency
relation of McKean--Vlasov form. In the operative uniform-prior case $\pi(\dd z'\mid z)=\rho(\dd z')$,
the first variation of \cref{eq:Frho} carries two distinct attention contributions, an outgoing
row term and an incoming source term, because the row normaliser $Z(z;\rho)$ depends on $\rho$ both
through the query at $z$ and through every other query that reads $z$ as a source,
\begin{equation}
\label{eq:firstvar}
\frac{\delta\Fcal}{\delta\rho}(y)
=\alpha\KL{q_y}{p}
-\tau\log Z(y)
-\tau\int\beta(z,y)\rho(\dd z)
+V(y),
\qquad
\beta(z,y)=\frac{e^{-\KL{q_z}{\Omega(z,y)q_y}/\tau}}{Z(z)},
\end{equation}
the term $-\tau\log Z(y)$ being the attention emitted by the query at $y$ over its sources and
$-\tau\int\beta(z,y)\rho(\dd z)$ the total attention received by $y$ from the population through
the incoming kernel $\mathcal{K}(\cdot,y)$; the directedness of the row coupling keeps these two terms
distinct, so no symmetric factor of $1/2$ appears. The Euler--Lagrange condition of
$\inf\{\KL{\rho}{\mu_1}:\Fcal[\rho]=f\}$ at the conjugate tilt $\bmeta$ then reads
$\log(\dd\rho/\dd\mu_1)=\mathrm{const}-\bmeta\delta\Fcal/\delta\rho$, that is the fixed-point
equation of the tilted single-site law
$\rho^{\star}(\dd y)\propto\mu_1(\dd y)\exp[-\bmeta\delta\Fcal/\delta\rho(y)]$ with $Z$ and
$\beta$ evaluated self-consistently at $\rho^{\star}$, the bounded-Lipschitz regularity this
characterisation presumes being supplied by the compact container. Identifying such a
stationary point with the continuum limit of the finite-population E-step iteration is a
separate, dynamical statement and does not follow from the microcanonical contraction
alone: the E-step as written is deterministic natural-gradient descent rather than a weakly
interacting diffusion, so the link requires a propagation-of-chaos or mean-field-limit
argument of the Dawson--G\"artner and Sznitman type
\citep{dawson1987large,sznitman1991topics}, whose bounded-Lipschitz hypotheses on the
interaction must then be checked on the container. The canonical alternative replaces the hard
constraint by a tilt: the Gibbs law $P_{\bmeta}\propto e^{-\bmeta\Fcal_{\mathrm{belief},N}}\mu_1^{\otimes N}$ over the same belief sector
has, by Sanov's theorem and Varadhan's lemma, the scaled-cumulant limit
$\tfrac1N\log\Zcal_N(\bmeta)\to-\varphi(\bmeta)$ with
$\varphi(\bmeta)=\inf_{\rho}\{\bmeta\Fcal[\rho]+\KL{\rho}{\mu_1}\}$, so $\varphi$ is the
negative of the $\bmeta$-scaled cumulant generating function (the negative pressure), not the free-energy density
$f$ itself. The microcanonical entropy density is recovered as the concave (infimum-form)
Legendre--Fenchel transform $s(f)=\inf_{\bmeta}\{\bmeta f-\varphi(\bmeta)\}+\text{const}$, the
transform direction that holds under the concavity (equivalence-of-ensembles) hypothesis; since
$\varphi$ is concave in $\bmeta$, the supremum-form transform does not recover $s$, and the two
routes should not be mixed \citep{ellis1985entropy,touchette2009large}. Working with the empirical law rather than
labelled agents quotients out permutations automatically, which resolves the
Gibbs-paradox bookkeeping and renders $\Smeta$ extensive without an inserted $1/N!$ when
the agents are exchangeable.

\section{The meta-temperature as natural-gradient noise}
\label{sec:temperature}

The Legendre conjugate $\bmeta=\partial\Smeta/\partial\Fcal_0$ defines the canonical
meta-ensemble $P[\omega]\propto e^{-\bmeta\Fcal_{\mathrm{belief}}^{\mathrm{red}}(\omega)}\mu_1^{\otimes N}$, a Gibbs
distribution over belief configurations weighted by their free energy, with partition
function $\Zcal_{\mathrm{meta}}=\int_{\Dcal^{N}} e^{-\bmeta\Fcal_{\mathrm{belief}}^{\mathrm{red}}}\mu_1^{\otimes N}(\dd\omega)$ and meta free energy
$\Phi=\langle\Fcal_{\mathrm{belief}}\rangle-\Tmeta\Smeta$, where the meta-temperature is the inverse
conjugate $\Tmeta=1/\bmeta$, so the inverse-temperature form $e^{-\bmeta\Fcal_{\mathrm{belief}}}$ and
the temperature form $e^{-\Fcal_{\mathrm{belief}}/\Tmeta}$ are the same weight. Because the container bounds the
free-energy density and $\Smeta$ need not be monotone in $\Fcal_0$, the canonical tilt $\bmeta=\partial\Smeta/\partial\Fcal_0$
is a signed parameter $\bmeta\in\R$, and the Riemannian-Langevin noise reading below applies only on the
positive branch $\bmeta=1/\Tmeta>0$; a negative-$\bmeta$ tilt remains a well-defined formal ensemble but
carries no diffusion temperature, since a noise amplitude cannot be negative. These weights are well defined only on the
container of \cref{sec:measure} where $\Zcal_{\mathrm{meta}}<\infty$. Rather than posit
this ensemble we identify it as the stationary law of a dynamics. Pure descent on
$\Fcal_{\mathrm{belief}}$ concentrates on the set of global minimisers, modulo the declared gauge quotient
and container constraint; it need not select a single configuration when minimisers are
degenerate or related by symmetry. Adding noise populates the ensemble, and on the curved
belief manifold the appropriate dynamics is the Riemannian Langevin equation, whose
It\^o form is
\begin{equation}
\label{eq:rlangevin}
\dd\omega^{i}
=\Bigl[-g^{ij}\partial_j\Fcal_{\mathrm{belief}}
+\Tmeta \tfrac{1}{\sqrt{\det g}} \partial_j\bigl((\sqrt{\det g}) g^{ij}\bigr)\Bigr]\dd t
+\sqrt{2\Tmeta}\bigl(\sqrt{g^{-1}}\bigr)^{i}{}_{k}\dd B^{k}.
\end{equation}
The first drift term $g^{ij}\partial_j\Fcal_{\mathrm{belief}}$ is the natural gradient, and the diffusion
covariance is $2\Tmeta g^{-1}$ per unit time, so $\Tmeta$ is the scalar temperature
multiplying the inverse metric rather than the full variance tensor. Subject to
confinement, non-explosiveness, and a normalisable partition function on the container,
the stationary law of \cref{eq:rlangevin} is the metric-weighted canonical ensemble
$P\propto e^{-\Fcal_{\mathrm{belief}}/\Tmeta}\sqrt{\det g}$, and this holds only when the metric-divergence
drift is retained; plain natural-gradient descent with isotropic noise does not sample it
\citep{girolami2011riemann,welling2011bayesian,xifara2014langevin}. This metric-weighted law equals
the declared canonical ensemble $e^{-\Fcal_{\mathrm{belief}}/\Tmeta}\mu_1^{\otimes N}$ of
\cref{eq:measure} only when the Riemannian volume reproduces the counting measure,
$\sqrt{\det g}\dd\omega=\dd\mathrm{Vol}_{\mathrm{FR}}(\bmu,\bSigma)\times\dd\mu_{\mathrm{Haar}}^{\Gcal}(U)$.
This is achieved by taking the product metric $g=g_{\mathrm{FR}}\oplus g_{\Gcal}$, where $g_{\mathrm{FR}}$
is the Fisher--Rao metric of \cref{sec:measure} and $g_{\Gcal}$ is the left-invariant metric on
$\Gcal=\GL^{+}(K)$ obtained by left-translating a fixed inner product on $\Lie(K)$: its Riemannian
volume is left Haar measure, which coincides with right Haar because $\GL^{+}(K)$ is unimodular, and
its pull-back to the exponential chart $U=e^{\bphi}$ is $\sqrt{\det g_{\Gcal}}\dd\bphi=J_{\exp}(\bphi)\dd\bphi$,
so the exp-chart Jacobian enters both the drift and the diffusion of \cref{eq:rlangevin} automatically
through $g$ and $\sqrt{\det g}$. Existence of this left-invariant metric is guaranteed; adopting it as
the Langevin metric is a modelling choice rather than a property of the E-step. On the hard container
of \cref{eq:container} the stationary law is exact only with reflecting boundary conditions on
$\partial\Dcal$, a formal prescription here, since the semialgebraic boundary carries nonsmooth strata
at repeated $\bSigma$ eigenvalues and frame singular-value multiplicities, where a well-posed reflected
diffusion requires a carefully specified Skorokhod problem or a smoothed domain. Replacing the hard
container by a smooth confining potential $V(\omega)\to\infty$ at the boundary removes the boundary
terms, but then samples the softened target $\propto e^{-(\Fcal_{\mathrm{belief}}+V)/\Tmeta}\sqrt{\det g}$,
the original hard-container law being recovered only in the infinite-barrier limit of $V$. The meta-temperature is therefore the
noise level injected into stochastic natural-gradient belief updating, an operational
quantity rather than an analogy. At $\Tmeta\to0$ the population concentrates on the
minimiser set; at finite $\Tmeta$ the meta-entropy is populated. The inverse
meta-temperature is the configuration-level counterpart of the precision that weights the
free energy in active-inference policy selection \citep{friston2017active}, lifted from
policies to whole belief configurations.

This places two Gibbs-weight temperatures and one relaxation scale in the theory, one for
each of its three timescales. The attention temperature $\tau=\kappa\sqrt{K}$ is the
soft-min temperature of the row log-partition $-\tau\log Z_i$ and governs within-configuration
inference on the fast E-step. The M-step learning rate is a relaxation and discretisation scale
for the slow adjustment of priors and models; a rate becomes a temperature only once an
accompanying injected-noise scale or fluctuation-dissipation relation is specified, as in
stochastic-gradient Langevin and constant-SGD samplers \citep{welling2011bayesian,mandt2017stochastic}.
The meta-temperature $\Tmeta$ is the genuine Gibbs weight of the ensemble over configurations,
the noise level of the Riemannian Langevin dynamics of \cref{eq:rlangevin}. The attention
temperature is already explicit in \cref{eq:F}; the meta-temperature is the new object of this
note.

\section{A solvable Gaussian response calculation}
\label{sec:phase}

A means-only Gaussian instance is solvable in closed form, and what it establishes is
narrow. It is a linear-response calculation, not a finite-coupling phase transition: the
order parameter vanishes at zero field for every finite prior strength, and the
susceptibility diverges only in the decoupling limit $\alpha\to0$. A genuine consensus
transition would require a compact or regulated model or a non-Gaussian covariance
sector, as discussed below.

\begin{proposition}[Gaussian meta-ensemble susceptibility]
\label{prop:gaussian}
Take isotropic fixed-covariance beliefs $q_i=\Ncal(\bmu_i,\sigma^2\bI)$, trivial gauge
$\Omegaij=\bI$, a normalised all-to-all weight $a_{ij}=1/N$ with coupling strength $J$,
and a quadratic prior field of strength $\alpha$ centred at the origin, so that up to a
$\sigma$-dependent constant
\begin{equation}
\label{eq:gaussianF}
\Fcal(\{\bmu\})=\frac{\alpha}{2}\sum_i\|\bmu_i\|^2
+\frac{J}{2}\sum_{i,j}a_{ij}\|\bmu_i-\bmu_j\|^2
=\frac{\alpha}{2}\sum_i\|\bmu_i\|^2+\frac{J}{2N}\sum_{i,j}\|\bmu_i-\bmu_j\|^2 .
\end{equation}
In the canonical meta-ensemble with the uniform field $b$ inserted as an energy-scale term inside the
free energy, $P\propto\exp[-\bmeta(\Fcal-b\cdot\sum_i\bmu_i)]$, conjugate to the magnetisation
$m=\frac1N\sum_i\bmu_i$, the mean-field
self-consistency reads $m=(2Jm+b)/(\alpha+2J)$, hence $m\alpha=b$. The zero-field
magnetisation is $m^{\star}=0$ and the susceptibility is
\begin{equation}
\chi=\frac{\partial m}{\partial b}=\frac{1}{\alpha},
\end{equation}
independent of the coupling $J$ and of $\bmeta$.
\end{proposition}

The $\bmeta$-independence is specific to the energy-scale convention above, in which the field and the free energy carry the common factor $\bmeta$ that cancels in the Gaussian mean. Under the source convention $P\propto\exp[-\bmeta\Fcal+b\cdot\sum_i\bmu_i]$, with $b$ entering the exponent unscaled, the same computation gives $m\alpha=b/\bmeta$ and $\chi=1/(\bmeta\alpha)$, the two conventions differing only by where the external field is scaled relative to $\bmeta$ \citep{kardar2007fields}.

\begin{proof}
Expanding the pairwise term,
$\sum_{i,j}\|\bmu_i-\bmu_j\|^2=2N\sum_i\|\bmu_i\|^2-2\|\sum_i\bmu_i\|^2$, so
$\Fcal=(\tfrac{\alpha}{2}+J)\sum_i\|\bmu_i\|^2-\tfrac{J}{N}\|\sum_i\bmu_i\|^2$. The
single-site energy in the mean field $m$ and external field $b$ is
$(\tfrac{\alpha}{2}+J)\|\bmu\|^2-(2Jm+b)\cdot\bmu$, a Gaussian with mean
$\langle\bmu\rangle=(2Jm+b)/(\alpha+2J)$. Imposing $\langle\bmu\rangle=m$ gives
$m(\alpha+2J)=2Jm+b$, that is $m\alpha=b$, whence $m^{\star}=0$ at $b=0$ and
$\chi=1/\alpha$.
\end{proof}

The susceptibility diverges as $\alpha\to0$: the prior strength is the control field, and
the means-only model has no spontaneous consensus at any finite $\alpha$. This
zero-mode susceptibility divergence is a linear-response signature, not an order--disorder
transition, and it is consistent with the gauge analysis of \cref{sec:gauge}, where
$\alpha$ is the strength of the explicit breaking of an otherwise exact covariance. A
genuine finite-coupling transition with a nonzero ordered magnetisation requires either a
spherical constraint $\sum_i\|\bmu_i\|^2=N$ in the manner of Berlin and Kac
\citep{berlin1952spherical} or the anharmonicity supplied by dynamical covariances. The
covariance sector is what makes the model non-Gaussian: a means-only theory is always
Gaussian, and the $\bSigma_i$ degrees of freedom would provide the nonlinearity required
for a sharp consensus transition. We therefore record the consensus transition itself as
a conjecture requiring a regulated or non-Gaussian model, and we do not claim it from
\cref{prop:gaussian}.

\section{The gauge sector: invariance, covariance, and redundancy}
\label{sec:gauge}

The order parameter is constrained by an exact invariance of the interaction, which we
establish below, and which the prior breaks. Three notions must be separated, because
only the last licenses gauge-fixing language. The first is the common-pushforward
invariance of the Gaussian divergence (\cref{lem:affine}). The second is the active
local covariance of the interaction under transforming all states and transports together
(\cref{prop:gaugeinv}). The third is true gauge redundancy, a transformation that leaves
the entire physical configuration unchanged and is to be quotiented in the partition
function. Whether the local transformations below are redundancies or active changes of
epistemic configuration depends on the gauge structure declared for the model. The
companion participatory manuscript treats the independent per-agent action
$\Omegaij\mapsto G_i\Omegaij G_j^{-1}$ explicitly and designates such per-agent frame
changes as active changes of epistemic configuration rather than redundancies; that
designation is a declared structural commitment, grounded in an edge-mode and
quantum-reference-frame analogy rather than derived from primitives, so we carry it as an
assumption and not as a settled result. We therefore prove the invariance and the active covariance, and we
keep the gauge-fixing and Goldstone reading explicitly conditional on the directions in
question being genuine redundancies.

\begin{lemma}[Affine invariance of the Gaussian divergence]
\label{lem:affine}
For invertible $G\in\GL(K)$ and $c\in\R^{K}$ let $T(x)=Gx+c$, so that
$T_{\#}\Ncal(\mu,\Sigma)=\Ncal(G\mu+c,G\Sigma G^{\top})$. Then for any two Gaussians
$q_1,q_2$,
\begin{equation}
\KL{T_{\#}q_1}{T_{\#}q_2}=\KL{q_1}{q_2}.
\end{equation}
\end{lemma}

\begin{proof}
The Gaussian divergence is
$\KL{q_1}{q_2}=\tfrac12[\Tr(\Sigma_2^{-1}\Sigma_1)+(\mu_2-\mu_1)^{\top}\Sigma_2^{-1}(\mu_2-\mu_1)-K+\log(\det\Sigma_2/\det\Sigma_1)]$.
Under $T$ the covariances map by congruence, $\Sigma_a\mapsto G\Sigma_a G^{\top}$ and
$\Sigma_2^{-1}\mapsto G^{-\top}\Sigma_2^{-1}G^{-1}$, so by cyclicity of the trace
$\Tr(\Sigma_2^{-1}\Sigma_1)$ is unchanged; the mean difference maps as
$\mu_2-\mu_1\mapsto G(\mu_2-\mu_1)$ and the quadratic form is invariant; and
$\det(G\Sigma_a G^{\top})=(\det G)^2\det\Sigma_a$, so the log-determinant ratio is
invariant. Each term is preserved.
\end{proof}

\begin{proposition}[Local covariance of the belief coupling]
\label{prop:gaugeinv}
Let the local transformation $g=(G_i)_{i=1}^{N}\in\GL^{+}(K)^{N}$ act by
$q_i\mapsto (G_i)_{\#}q_i$ and $e^{\bphi_i}\mapsto G_i e^{\bphi_i}$. This is a
pre-regulator algebraic identity on the ambient group: the underlying common-pushforward
invariance of \cref{lem:affine} holds for all of $\GL(K)$, and the proposition is a
covariance of the unregulated interaction, not of the regulated counted state space.
The action does not preserve the counted gauge chart of \cref{eq:nesting}: a noncompact
$G_i$ carries $G_ie^{\bphi_i}$ outside $\{e^{\bphi}:\|\bphi\|_2\le\pi-\eta\}$, a $G_i$ of
negative determinant carries it outside $\GL^{+}(K)$, and $G_ie^{\bphi_i}$ need not be an
exponential at all \citep{culver1966existence}; the resulting failure to preserve the
container is the explicit regulator breaking discussed after \cref{prop:breaking}. Then the transport
maps as $\Omegaij\mapsto G_i\Omegaij G_j^{-1}$, the transported source belief satisfies
$(G_i\Omegaij G_j^{-1})_{\#}(G_j)_{\#}q_j=(G_i)_{\#}(\Omegaij)_{\#}q_j$, and each summand
$\KL{q_i}{\Omegaij q_j}$ is invariant. Consequently the belief-coupling functional
$\sum_{i,j}\betaij \KL{q_i}{\Omegaij q_j}$ is invariant, the softmax weights $\betaij$
being functions of the invariant divergences.
\end{proposition}

\begin{proof}
Transport acts on a Gaussian by the congruence
$\Omegaij\cdot\Ncal(\mu,\Sigma)=\Ncal(\Omegaij\mu,\Omegaij\Sigma\Omegaij^{\top})$, that is
by pushforward under the linear map $\Omegaij$. The transformed transport is
$G_i\Omegaij G_j^{-1}$; composing its pushforward with that of $G_j$ gives pushforward
under $(G_i\Omegaij G_j^{-1})G_j=G_i\Omegaij$, hence
$(G_i\Omegaij)_{\#}q_j=(G_i)_{\#}(\Omegaij q_j)$. The target transforms as
$q_i\mapsto(G_i)_{\#}q_i$. Both arguments of the divergence acquire the common map $G_i$,
so by \cref{lem:affine} the summand is unchanged.
\end{proof}

The covariance is exact and algebraic: \cref{prop:gaugeinv} follows from \cref{lem:affine}
in two lines, since both arguments of each summand acquire the common map $G_i$, while the
self-coupling does not co-transform and so is not invariant.

\begin{proposition}[The prior is the symmetry-breaking field]
\label{prop:breaking}
The self-coupling $\alpha\KL{q_i}{p_i}$ with fixed prior $p_i$ is invariant under
$q_i\mapsto(G_i)_{\#}q_i$ only if the prior co-transforms, $p_i\mapsto(G_i)_{\#}p_i$.
With the priors held fixed as M-step references, the local covariance of the
belief-sector interaction $\Fcal_{\mathrm{belief}}$ is explicitly broken by two terms: the
prior-anchoring term $\alpha\sum_i\KL{q_i}{p_i}$, whose strength is set by $\alpha$, and the
observation term $-\E_q[\log p(o\mid x)]$, which is not invariant under $q_i\mapsto(G_i)_{\#}q_i$
unless the channel co-transforms, since the data $o$ and the likelihood $p(o\mid x)$ do not
transform with the frame. The hyper-prior term $\lambda_h\sum_i\KL{s_i}{h}$ does not appear in
$\Fcal_{\mathrm{belief}}$ and so plays no role in the breaking of the counted belief covariance.
The $s_i/h$ and model-coupling $\gammaij$ sector of the full $\Fcal$ acts on the frame variables
as well, and its symmetry analysis is separate; we do not undertake it here.
\end{proposition}

The consequences depend on whether the transformations of \cref{prop:gaugeinv} are
genuine gauge redundancies. If they are, then in the limit $\alpha\to0$, and for a gauge-invariant or absent observation
term so that the likelihood imposes no residual breaking, the belief-sector redundancy is
exact, the configuration measure has flat directions along the gauge orbits, and
$\Wmeta(\Fcal_0)$ diverges through the noncompact gauge volume unless one quotients by the
gauge group; this quotient is the Faddeev--Popov gauge
fixing of the partition function familiar from gauge field theory
\citep{faddeev1967feynman}, and for $\alpha>0$ the broken directions become
pseudo-Goldstone soft modes whose contribution grows smoothly as $\alpha\to0$. At the fixed compact
container $\Dcal$ this contribution saturates at the finite soft-direction volume rather than diverging;
the unbounded gauge-volume divergence appears only once the regulator is removed, taking $\alpha\to0$ first
and the container scale to infinity afterward. The growth is consistent with the susceptibility
$\chi=1/\alpha$ of \cref{prop:gaussian}. If instead
the transformations are active changes of epistemic configuration, the directions are
physical and must not be quotiented; counting them is then correct but the
Goldstone/gauge-fixing language is inappropriate, and the $\alpha\to0$ divergence is
read as a genuine soft-mode susceptibility rather than a gauge volume. In either reading
the noncompactness of $\R^{K}$, $\SPD(K)$ and $\Gcal$ means that $\Wmeta(\Fcal_0)$ is
finite only relative to a compact container. Four distinct things must be kept apart. The
first is the invariance of the unregulated interaction: each summand $\KL{q_i}{\Omegaij q_j}$
is invariant under the common pushforward of \cref{prop:gaugeinv}, an algebraic identity on the
ambient group. The second concerns the counting measure: the ambient density $\dd\mu_1$
of \cref{eq:measure} is left-invariant in form before restriction, since the ambient Haar density is
invariant under the left translation $U\mapsto G U$ and the affine-invariant Fisher--Rao volume is
preserved by the congruence $\bSigma\mapsto G\bSigma G^{\top}$ \citep{pennec2006riemannian}. After
restriction to the counted chart the relevant measure is the restricted Haar density, and a generic
left translation no longer maps the counted domain to itself, so the local action is not a symmetry of
the measured domain; this loss is the container breaking recorded as the fourth mechanism below.
The third is explicit breaking within $\Fcal_{\mathrm{belief}}$ by the prior anchoring, of strength $\alpha$, and by the observation likelihood, as in \cref{prop:breaking}. The fourth is explicit breaking
by the finite container: a noncompact $G_i$ can push $\bmu_i$ past $\|\bmu_i\|\le R$, move the
eigenvalues of $\bSigma_i$ outside the two-sided bounds, or carry $G_ie^{\bphi_i}$ outside the
principal-log chart. A regulator is expected to break a noncompact symmetry, so the container
breaking is a feature of the construction rather than a defect; the gauge-volume and
pseudo-Goldstone reading above is attached to the explicit prior breaking of the third kind, with
the fourth kind always present once $\Dcal$ is imposed. The bare precision bound
$\tr\bSigma_i^{-1}\le c$ does not suffice: it constrains each $\bSigma_i$ from below only
and leaves $\KL{q_i}{p_i}\to+\infty$ as $\bSigma_i\to\infty$. We therefore take the
two-sided set
\begin{equation}
\label{eq:container}
\Dcal=\bigl\{c_1\bI\preceq\bSigma^{-1}\preceq c_2\bI,\ \ \|\bmu\|\le R,\ \ \|\bphi\|_2\le\pi-\eta\bigr\},
\qquad 0<c_1\le c_2<\infty,\ \ R<\infty,\ \ 0<\eta<\pi,
\end{equation}
the single-agent container, with $\Dcal^{N}$ the product container used in the microcanonical
integral and the reference law; equivalently $(1/c_2)\bI\preceq\bSigma\preceq(1/c_1)\bI$, which
bounds the covariance from both sides; the closed spectral bound $\|\bphi\|_2\le\pi-\eta$ confines
the frame $U=e^{\bphi}$ to the counted gauge chart $\{e^{\bphi}:\|\bphi\|_2\le\pi-\eta\}$ of
\cref{eq:nesting}, a compact subset of the exponential image, itself a proper subset of
$\GL^{+}(K)$ \citep{culver1966existence}, lying strictly inside the open injective exponential chart
$\{\|\bphi\|_2<\pi\}$ on which the principal logarithm inverts $\exp$ with non-vanishing
Jacobian \citep{higham2008functions}. On
$\Dcal$ the affine-invariant Fisher--Rao volume and the Haar mass of the compact gauge
image are finite, so the reference $\mu_1$ restricted to $\Dcal$ normalises to a
probability law, the analogue of the finite volume an ideal gas requires; the meta-entropy
is well defined once such a constraint is declared, and this should precede any
theorem-like use of a finite $\Wmeta(\Fcal_0)$.

\section{Holonomy, frustration, and residual meta-entropy}
\label{sec:frustration}

The transport in \cref{eq:F} is vertex-frame, $\Omegaij=g_i g_j^{-1}$ with
$g_i=e^{\bphi_i}$. Around any closed loop the product telescopes,
\begin{equation}
\Omega_{i_1 i_2}\Omega_{i_2 i_3}\cdots\Omega_{i_n i_1}
=g_{i_1}g_{i_2}^{-1}g_{i_2}g_{i_3}^{-1}\cdots g_{i_n}g_{i_1}^{-1}=\bI,
\end{equation}
so the holonomy is identically trivial in the regime the functional actually defines, by the
exact telescoping of the displayed product, matching the vanishing-holonomy result established
for vertex-frame transport in the companion attention manuscript. No
residual holonomy entropy can arise in this regime. A nontrivial holonomy requires an
edge-local connection that does not telescope, for example
$\Omegaij=e^{\bphi_i}e^{\delta_{ij}G}e^{-\bphi_j}$ with an edge generator $\delta_{ij}G$,
which inserts a genuine loop curvature. The conjecture below is therefore a statement
about this distinct, edge-relaxed regime, not about \cref{eq:F} as written.

Even in the edge-local regime, nonzero holonomy does not by itself produce a residual
entropy. A holonomy group can have fixed vectors, and a frustrated constraint system can
still possess a unique ground state; what is required is an extensive ground-state
degeneracy, a near-ground manifold whose dimension grows with $N$. In the continuous
belief sector frustration can act in the opposite direction: a uniform twist gaps the
consensus soft mode, raising the lowest eigenvalue of the coupling operator and lowering,
rather than raising, the configurational entropy. A positive residual entropy density is
thus a property of the ground-state degeneracy, not of holonomy alone, and typically
needs a compact or regulated sector to be realised. We state the conditional consequence
as a conjecture.

\begin{conjecture}[Holonomy-induced residual meta-entropy, edge-local regime]
\label{conj:frustration}
Adopt an edge-local connection with holonomy $\holonomy(\gamma)$ nontrivial on a set of
independent loops $\gamma$ of extensive cardinality, and a compact or regulated belief
sector. Suppose the induced alignment constraints are mutually incompatible in the sense
that the near-ground set is an extensive flat manifold, or a discretely resolved family under the
declared reference measure, whose log-volume grows linearly with $N$, rather than a collection of
isolated minima with merely soft harmonic directions. Then the
consensus order parameter is frustrated and the meta-entropy retains a strictly positive
density at zero meta-temperature in the thermodynamic order of limits,
$\lim_{\Tmeta\to0}\lim_{N\to\infty} N^{-1}\Smeta>0$ (population limit first, then
$\Tmeta\to0$), the analogue of the residual entropy of ice \citep{pauling1935structure} and
of spin-glass
ground-state degeneracy \citep{mezard1987spin}. For a flat connection, or whenever the
ground state is unique up to the declared quotient, the residual density vanishes.
\end{conjecture}

The conjecture is falsifiable by simulation in the edge-local regime. One arranges agents
on a frustrated lattice with an extensive number of independent plaquettes, compares a
flat connection against one carrying nonzero plaquette holonomy, samples the canonical
meta-ensemble with the Riemannian Langevin dynamics of \cref{eq:rlangevin} on a regulated
sector, and estimates the residual entropy density by thermodynamic integration of
$\langle\Fcal_{\mathrm{belief}}\rangle$ over $\bmeta$. Because
$\partial_{\bmeta}\log\Zcal_{\mathrm{meta}}=-\langle\Fcal_{\mathrm{belief}}\rangle$, this integration is
anchored at $\bmeta=0$, where $\log\Zcal_{\mathrm{meta}}(0)=0$ because the reference
$\mu_1^{\otimes N}$ normalises to a probability law on the container of \cref{eq:container}, so no
free integration constant remains and the entropy is fixed relative to that declared reference,
the same reference carried as the $+\text{const}$ of \cref{eq:contraction}. The quantity the test
estimates is the canonical entropy density, not the log-partition alone: at inverse meta-temperature
$\bmeta$ the canonical meta-entropy is
\begin{equation}
\label{eq:scan}
S_{\mathrm{can}}(\bmeta)=\log\Zcal_{\mathrm{meta}}(\bmeta)+\bmeta\langle\Fcal_{\mathrm{belief}}\rangle_{\bmeta},
\end{equation}
both terms accumulated along the same integration, and the residual density is its zero-temperature
limit $s_0=\lim_{\bmeta\to\infty}N^{-1}[\log\Zcal_{\mathrm{meta}}(\bmeta)+\bmeta\Fcal_{\mathrm{belief},0}]$
when the free-energy density tends to the ground value $f_0=\Fcal_{\mathrm{belief},0}/N$. This canonical
$s_0$ coincides with the microcanonical ground-shell entropy density
$\lim_{N\to\infty}N^{-1}\Smeta(\Fcal_{\mathrm{belief},0})$ of \cref{eq:Smicro} only where the
Legendre--Fenchel equivalence of ensembles of \cref{sec:meta} holds, which can fail at the ground
value if concavity is lost. Because the belief sector is continuous, the test must also separate the
harmonic low-temperature contribution of the soft directions around each minimum, which scales as a
$\log\Tmeta$ term per extensive normal mode under the declared reference, from the residual flat or
discretely resolved degeneracy that survives the $\Tmeta\to0$ limit as a finite log-volume; only the
latter enters the conjectured positive density. The holonomic-minus-flat plateau difference targeted
by the test is in any case reference independent.
A positive plateau in the holonomic case and its absence in the flat case would support
\cref{conj:frustration}; absence in both, or a plateau that vanishes as the regulator is removed,
would refute it.

\section{Scope and open problems}
\label{sec:discussion}

The construction supports a qualified answer to the original question. There is a
candidate configurational meta-entropy $\Smeta=\log\Wmeta(\Fcal_0)$, counted on the
belief sector against the Fisher--Rao volume and on the gauge sector against a separately
specified group measure with a quotient or regulator. It is rendered extensive in the
large-population limit through the row log-partition $-\tau\log Z_i$ of \cref{eq:rowfree},
with the $1/N$ Kac normalisation retained inside the logarithm; its conjugate temperature
is the noise level of stochastic natural-gradient updating; and the structure that
distinguishes it from ordinary thermodynamics is the gauge sector, whose soft modes are
pseudo-Goldstone directions only when they are genuine redundancies, and whose frustration
is governed by the curvature of an edge-local connection.

We separate three confidence levels. Exact and verified in this note are the affine
invariance of the Gaussian divergence (\cref{lem:affine}), the local covariance of the
belief coupling and its breaking by the fixed prior and the fixed observation channel
(\cref{prop:gaugeinv,prop:breaking}), the susceptibility
$\chi=1/\alpha$ of the Gaussian response calculation (\cref{prop:gaussian}), and the
telescoping of vertex-frame holonomy to the identity. Standard and imported, not
re-derived here, are the Kac criterion for extensivity, the Chentsov uniqueness of the
Fisher--Rao metric, the Sanov and contraction apparatus of \cref{sec:meta} together with
the McKean--Vlasov stationary form of its constrained minimiser, the Riemannian-Langevin stationary law with its
metric-divergence drift under confinement, and the active-inference reading of $\bmeta$.
Conjectural, and explicitly conditional on an edge-local connection and on extensive
ground-state degeneracy, is the holonomy-frustration mechanism of
\cref{conj:frustration}; conjectural likewise is the consensus phase transition, which
the Gaussian model does not deliver and which requires a regulated or non-Gaussian
covariance sector.

Several problems remain open. The container constraint that regularises the noncompact
gauge and scale directions should be chosen so that the meta-entropy is a smooth function
of $\Fcal_0$; the natural candidate, bounded average precision, has a clear
information-geometric meaning but its effect on the response calculation has not been
computed. The Sanov rate of \cref{eq:contraction} closes only partially for the Gaussian
family (\cref{rem:gaussianrate}): the rate functional reduces to a negative Fisher--Rao
entropy and the affine-invariant geometry of $\SPD(K)$ supplies the closed geodesic of the
covariance sector, but the constrained minimiser is elementary only when the covariance
sector is frozen, and the covariance-active minimiser remains an open variational problem;
the equivalence of the microcanonical and canonical routes should also be checked where
concavity can fail. The status of the local transformations of
\cref{prop:gaugeinv} as redundancies or as active changes should be settled against the
operative gauge actions of the wider theory, since the gauge-fixing language of
\cref{sec:gauge} depends on it. The coupling between the meta-temperature and the
attention temperature $\tau$ invites a Maxwell-relation analysis, since both are
temperatures in the same hierarchy. Finally, \cref{conj:frustration} should be tested in
the edge-local regime on a frustrated lattice with a regulated belief sector, which is
the natural next step.

\begin{thebibliography}{99}

\bibitem{amari2000methods}
S.~Amari and H.~Nagaoka.
\newblock \emph{Methods of Information Geometry}.
\newblock American Mathematical Society and Oxford University Press, 2000.

\bibitem{berlin1952spherical}
T.~H. Berlin and M.~Kac.
\newblock The spherical model of a ferromagnet.
\newblock \emph{Physical Review}, 86(6):821--835, 1952.

\bibitem{chentsov1982statistical}
N.~N. Chentsov.
\newblock \emph{Statistical Decision Rules and Optimal Inference}.
\newblock Translations of Mathematical Monographs, vol.~53, American Mathematical
  Society, 1982.

\bibitem{cover2006elements}
T.~M. Cover and J.~A. Thomas.
\newblock \emph{Elements of Information Theory}.
\newblock Wiley, 2nd edition, 2006.

\bibitem{csiszar1975iprojection}
I.~Csisz\'ar.
\newblock $I$-divergence geometry of probability distributions and minimization problems.
\newblock \emph{The Annals of Probability}, 3(1):146--158, 1975.

\bibitem{kardar2007fields}
M.~Kardar.
\newblock \emph{Statistical Physics of Fields}.
\newblock Cambridge University Press, 2007.

\bibitem{comets1989large}
F.~Comets.
\newblock Large deviation estimates for a conditional probability distribution.
  Applications to random interaction Gibbs measures.
\newblock \emph{Probability Theory and Related Fields}, 80(3):407--432, 1989.

\bibitem{culver1966existence}
W.~J. Culver.
\newblock On the existence and uniqueness of the real logarithm of a matrix.
\newblock \emph{Proceedings of the American Mathematical Society}, 17(5):1146--1151, 1966.

\bibitem{dawson1987large}
D.~A. Dawson and J.~G\"artner.
\newblock Large deviations from the McKean--Vlasov limit for weakly interacting
  diffusions.
\newblock \emph{Stochastics}, 20(4):247--308, 1987.

\bibitem{dembo1998large}
A.~Dembo and O.~Zeitouni.
\newblock \emph{Large Deviations Techniques and Applications}.
\newblock Springer, 2nd edition, 1998.

\bibitem{ellis1985entropy}
R.~S. Ellis.
\newblock \emph{Entropy, Large Deviations, and Statistical Mechanics}.
\newblock Springer, 1985.

\bibitem{faddeev1967feynman}
L.~D. Faddeev and V.~N. Popov.
\newblock Feynman diagrams for the Yang--Mills field.
\newblock \emph{Physics Letters B}, 25(1):29--30, 1967.

\bibitem{friston2010free}
K.~Friston.
\newblock The free-energy principle: a unified brain theory?
\newblock \emph{Nature Reviews Neuroscience}, 11(2):127--138, 2010.

\bibitem{friston2017active}
K.~Friston, T.~FitzGerald, F.~Rigoli, P.~Schwartenbeck, and G.~Pezzulo.
\newblock Active inference: a process theory.
\newblock \emph{Neural Computation}, 29(1):1--49, 2017.

\bibitem{girolami2011riemann}
M.~Girolami and B.~Calderhead.
\newblock Riemann manifold Langevin and Hamiltonian Monte Carlo methods.
\newblock \emph{Journal of the Royal Statistical Society: Series B}, 73(2):123--214, 2011.

\bibitem{xifara2014langevin}
T.~Xifara, C.~Sherlock, S.~Livingstone, S.~Byrne, and M.~Girolami.
\newblock Langevin diffusions and the Metropolis-adjusted Langevin algorithm.
\newblock \emph{Statistics \& Probability Letters}, 91:14--19, 2014.

\bibitem{goldstein1980classical}
H.~Goldstein.
\newblock \emph{Classical Mechanics}.
\newblock Addison-Wesley, 2nd edition, 1980.

\bibitem{hall2015lie}
B.~C. Hall.
\newblock \emph{Lie Groups, Lie Algebras, and Representations: An Elementary Introduction}.
\newblock Graduate Texts in Mathematics, vol.~222, Springer, 2nd edition, 2015.

\bibitem{higham2008functions}
N.~J. Higham.
\newblock \emph{Functions of Matrices: Theory and Computation}.
\newblock Society for Industrial and Applied Mathematics, 2008.

\bibitem{kac1963vanderwaals}
M.~Kac, G.~E. Uhlenbeck, and P.~C. Hemmer.
\newblock On the van der Waals theory of the vapor--liquid equilibrium. I.
\newblock \emph{Journal of Mathematical Physics}, 4(2):216--228, 1963.

\bibitem{mandt2017stochastic}
S.~Mandt, M.~D. Hoffman, and D.~M. Blei.
\newblock Stochastic gradient descent as approximate Bayesian inference.
\newblock \emph{Journal of Machine Learning Research}, 18(134):1--35, 2017.

\bibitem{mezard1987spin}
M.~M\'ezard, G.~Parisi, and M.~A. Virasoro.
\newblock \emph{Spin Glass Theory and Beyond}.
\newblock World Scientific, 1987.

\bibitem{nielsen2020elementary}
F.~Nielsen.
\newblock An elementary introduction to information geometry.
\newblock \emph{Entropy}, 22(10):1100, 2020.

\bibitem{pennec2006riemannian}
X.~Pennec, P.~Fillard, and N.~Ayache.
\newblock A Riemannian framework for tensor computing.
\newblock \emph{International Journal of Computer Vision}, 66(1):41--66, 2006.

\bibitem{pinele2020fisher}
J.~Pinele, J.~E. Strapasson, and S.~I.~R. Costa.
\newblock The Fisher--Rao distance between multivariate normal distributions: special cases, bounds and applications.
\newblock \emph{Entropy}, 22(4):404, 2020.

\bibitem{skovgaard1984riemannian}
L.~T. Skovgaard.
\newblock A Riemannian geometry of the multivariate normal model.
\newblock \emph{Scandinavian Journal of Statistics}, 11(4):211--223, 1984.

\bibitem{pauling1935structure}
L.~Pauling.
\newblock The structure and entropy of ice and of other crystals with some randomness of
  atomic arrangement.
\newblock \emph{Journal of the American Chemical Society}, 57(12):2680--2684, 1935.

\bibitem{sznitman1991topics}
A.-S. Sznitman.
\newblock Topics in propagation of chaos.
\newblock In \emph{\'Ecole d'\'Et\'e de Probabilit\'es de Saint-Flour XIX--1989},
  Lecture Notes in Mathematics, vol.~1464, pp.~165--251. Springer, 1991.

\bibitem{touchette2009large}
H.~Touchette.
\newblock The large deviation approach to statistical mechanics.
\newblock \emph{Physics Reports}, 478(1--3):1--69, 2009.

\bibitem{vaswani2017attention}
A.~Vaswani, N.~Shazeer, N.~Parmar, J.~Uszkoreit, L.~Jones, A.~N. Gomez, L.~Kaiser, and
  I.~Polosukhin.
\newblock Attention is all you need.
\newblock In \emph{Advances in Neural Information Processing Systems 30}, 2017.

\bibitem{welling2011bayesian}
M.~Welling and Y.~W. Teh.
\newblock Bayesian learning via stochastic gradient Langevin dynamics.
\newblock In \emph{Proceedings of the 28th International Conference on Machine Learning},
  pp.~681--688, 2011.

\end{thebibliography}

\end{document}
